\section{Results}\label{sec:results}
We report the \emph{primary} perceived-stress analyses in the KNHANES development
cohort, then the \emph{supporting} NHANES analyses. KNHANES is the primary test
because, unlike NHANES, it carries a general perceived-stress item (the NHANES
PHQ-9 proxy is a DAG \emph{mediator}). All estimates are design-based
(survey-weighted, Taylor-linearized) associations, not causal effects, and every
value below is reproduced by the committed \texttt{outputs/knhanes\_results.json}
and \texttt{nhanes\_results.json}.

\subsection{Primary results (KNHANES, adults 40--69)}
Pooling the otologic-examination cycles 2010--2012 (KNHANES V; analytic domain
$n=10{,}467$; complete-case models retained 9{,}142--9{,}532 by outcome),
survey-weighted prevalence was 21.1\% (95\% CI 19.9--22.3) for any tinnitus, 7.0\%
(6.3--7.7) for bothersome tinnitus, 15.4\% (14.5--16.4) for better-ear
speech-frequency hearing loss ($>$25 dB HL), and 24.3\% (23.2--25.3) for high
perceived stress.
\IfFileExists{tables/table_results_knhanes.tex}{\begin{table}[htbp]
\centering
\caption{\textbf{Real KNHANES results} (Korean adults 40--69). Primary development cohort with the perceived-stress exposure. Prevalence is design-based (survey-weighted, Taylor-linearized); associations are mutually adjusted design-based odds ratios and are \emph{not} causal. Generated by \texttt{code/src/knhanes\_analysis.py}. $^{*}p<0.05$, $^{**}p<0.01$, $^{***}p<0.001$.}
\label{tab:results_knhanes}
\small
\begin{tabular}{lcc}
\toprule
 & Tinnitus & Hearing loss ($>$25 dB HL) \\
\midrule
Prevalence (95\% CI) & 21.1\% (19.9--22.3) & 15.4\% (14.5--16.4) \\
\midrule
\multicolumn{3}{l}{\emph{Adjusted associations, OR (95\% CI)}} \\
\quad Perceived stress (high) & 1.42 (1.26--1.60)*** & 1.18 (0.99--1.41) \\
\quad Occupational noise & 1.37 (1.15--1.63)*** & 1.72 (1.44--2.06)*** \\
\quad Age (per year) & 1.03 (1.02--1.04)*** & 1.13 (1.12--1.14)*** \\
\quad Female sex & 1.08 (0.96--1.22) & 0.47 (0.40--0.55)*** \\
\bottomrule
\end{tabular}
\end{table}
}{}

\paragraph{Primary association (H1).} In the minimal-sufficient model (perceived
stress, occupational noise, age, sex), high perceived stress was associated with
\emph{tinnitus} (OR 1.42, 1.26--1.60) but only weakly and non-significantly with
\emph{better-ear hearing loss} (OR 1.18, 0.99--1.41), where occupational noise (OR
1.72, 1.44--2.06) and age (OR 1.13 per year) dominated. This is the prespecified
symptom-versus-threshold dissociation (H1): perceived stress tracks the tinnitus
\emph{symptom} ($p<0.001$) but attenuates against the audiometric \emph{threshold}
($p=0.07$) once occupational noise and age are accounted for.

\paragraph{Sensitivity analyses.} The stress--tinnitus association survived
adjustment for the mediator depressed mood (OR 1.29, 1.13--1.47) and for
audiometric hearing loss itself (OR 1.44, 1.27--1.64), so it is not an artifact of
under-adjustment for the strongest tinnitus correlate
(Table~\ref{tab:results_knhanes_extended}). It persisted for \emph{non-bothersome}
tinnitus versus none (OR 1.24, 1.05--1.45) and was strongest for bothersome
tinnitus (OR 1.79, 1.46--2.21); across seven secondary contrasts we control the
false-discovery rate (Benjamini--Hochberg), and these remain significant
($q<0.02$). The dissociation is \emph{relative, not absolute}: perceived stress
was weakly associated with the \emph{worse-ear} threshold (OR 1.15, 1.00--1.31;
does not survive FDR, $q\approx0.07$) but not with the high-frequency definition
(OR 1.07, 0.93--1.23). Because the KNHANES examination provides air-conduction
thresholds and tympanometry but not bone conduction, we excluded middle-ear
pathology using normal bilateral tympanometry: the stress--threshold association
then \emph{strengthened} to significance (OR 1.29, 1.07--1.55) while the tinnitus
association was essentially unchanged (OR 1.42, 1.25--1.62), indicating the
full-sample null was partly a conductive-contamination artifact and that the
contrast is one of \emph{degree}.
\IfFileExists{tables/table_extended_knhanes.tex}{\begin{table}[htbp]
\centering
\caption{\textbf{Extended and sensitivity models (KNHANES, adults 40--69).}
Perceived-stress odds ratio (95\% CI) across prespecified specifications,
isolating the robustness of the primary exposure. Rows add the mediator
(depressed mood) or audiometric hearing loss to the minimal-sufficient set
(perceived stress, occupational noise, age, sex), and vary the hearing-loss
outcome definition (better-ear, worse-ear, high-frequency). All are design-based
associations, not causal effects; mediator-adjusted rows are interpreted as
mediator- (not confounder-) adjusted. Generated by
\texttt{code/src/knhanes\_analysis.py}. Cells marked \emph{[pending data run]}
populate automatically when the pipeline is executed on the KDCA-approved
microdata with \texttt{-{}-latex-extended}. $^{*}p<0.05$, $^{**}p<0.01$,
$^{***}p<0.001$.}
\label{tab:results_knhanes_extended}
\small
\begin{tabular}{lcc}
\toprule
Model specification & Perceived stress, OR (95\% CI) & $N$ \\
\midrule
Tinnitus --- minimal-sufficient (primary) & 1.42 (1.26--1.60)*** & 9504 \\
\quad + depressed mood (mediator-adjusted) & \textit{[pending data run]} & -- \\
\quad + audiometric hearing loss & \textit{[pending data run]} & -- \\
Hearing loss, better ear --- minimal-sufficient & 1.18 (0.99--1.41) & 9176 \\
\quad + depressed mood (mediator-adjusted) & \textit{[pending data run]} & -- \\
Hearing loss, worse ear --- minimal-sufficient & \textit{[pending data run]} & -- \\
Hearing loss, high-freq (3/4/6\,kHz) --- minimal-sufficient & \textit{[pending data run]} & -- \\
Bothersome tinnitus --- minimal-sufficient & \textit{[pending data run]} & -- \\
Non-bothersome tinnitus vs none (reverse-causation sensitivity) & \textit{[pending data run]} & -- \\
\bottomrule
\end{tabular}
\end{table}
}{}

The employed-restricted (healthy-worker) contrast, the KNHANES$\rightarrow$NHANES
external validation (H2), and subgroup fairness are prospective components that
follow the prespecified plan (Methods; Supplementary Methods~S1) and carry no
results here; the open-LLM extraction and referral-safety layer are developed,
with results, in the companion SAFE-EAR paper.

\subsection{Supporting results (NHANES): directional consistency}
For U.S. adults aged 40--69 (pooling 2011--2012, 2015--2016, 2017--2018;
Table~\ref{tab:results}), survey-weighted prevalence was 28.5\% (26.5--30.4) for
any tinnitus, 12.0\% (10.4--13.6) for better-ear hearing loss, and 8.4\% (7.5--9.4)
for a positive PHQ-9 screen. In mutually adjusted design-based models,
\emph{tinnitus} was associated with occupational noise (OR 1.59, 1.32--1.92),
depressed mood (OR 1.61, 1.23--2.11), and female sex (OR 1.61, 1.35--1.91) but not
age; \emph{hearing loss} was dominated by age (OR 1.11 per year, 1.09--1.13) and
depressed mood (OR 2.29, 1.56--3.35), inversely by female sex (OR 0.42,
0.29--0.61), and only weakly by noise (OR 1.22, 0.94--1.58). Because the PHQ-9
proxy is a DAG mediator and NHANES lacks the perceived-stress exposure, these
mediator-level associations reproduce the symptom-versus-threshold pattern in a
different population but do not substitute for the primary KNHANES analysis.
\IfFileExists{tables/table_results.tex}{\begin{table}[htbp]
\centering
\caption{Real NHANES survey-weighted results (adults 40--69). Estimates are design-based (Taylor-linearized); associations are odds ratios from design-weighted logistic regression and are not causal. Generated by \texttt{code/src/nhanes\_analysis.py}.}
\label{tab:results}
\small
\begin{tabular}{lll}
\toprule
Quantity & Estimate (95\% CI) & $N$ \\
\midrule
Tinnitus prevalence & 28.5\% (26.5--30.4) & 5354 \\
Hearing loss ($>$25 dB HL) prevalence & 12.0\% (10.4--13.6) & 4808 \\
Depression (PHQ-9$\geq$10) prevalence & 8.4\% (7.5--9.4) & 7293 \\
Tinnitus: OR depressed mood & 1.61 (1.23--2.11) & 4603 \\
Tinnitus: OR occupational noise & 1.59 (1.32--1.92) & 4603 \\
Hearing loss: OR depressed mood & 2.29 (1.56--3.35) & 4463 \\
Hearing loss: OR occupational noise & 1.22 (0.94--1.58) & 4463 \\
\bottomrule
\end{tabular}
\end{table}
}{}
